\documentclass[aspectratio=169,11pt]{beamer}
\usetheme{Madrid}
\usecolortheme{dolphin}
\setbeamertemplate{navigation symbols}{}
\setbeamertemplate{caption}[numbered]

\usepackage[utf8]{inputenc}
\usepackage[T1]{fontenc}
\usepackage[spanish,es-noshorthands]{babel}
\usepackage{amsmath,amssymb,amsthm,mathtools}
\usepackage{tikz}
\usepackage{pgfplots}
\pgfplotsset{compat=1.18}
\usepackage{booktabs}
\usepackage{array}

\title[Prob. y Estadística]{Clase 8: Variable aleatoria continua, distribución Uniforme y distribución Normal}
\subtitle{Unidad 2: Variables aleatorias}
\institute[UNS]{Universidad Nacional del Sur \\ Departamento de Matemática}
\author{Probabilidad y Estadística (5765)}
\date{}

\begin{document}

\begin{frame}
\titlepage
\end{frame}

\begin{frame}{Contenidos}
\tableofcontents
\end{frame}

\section{Variable aleatoria continua}

\begin{frame}{De lo discreto a lo continuo}
\begin{block}{Motivación}
Muchas variables de interés (tiempo, longitud, peso, temperatura) toman \textbf{cualquier valor} en un intervalo real: no tiene sentido hablar de $P(X=x)$ como en el caso discreto, pues ese valor es $0$ para cada $x$.
\end{block}

\begin{definition}[Variable aleatoria continua]
$X$ es una v.a. \textbf{continua} si existe una función $f:\mathbb{R}\to[0,\infty)$, llamada \textbf{función de densidad de probabilidad}, tal que
$$F(x) = P(X\le x) = \int_{-\infty}^x f(t)\,dt \qquad \text{para todo } x\in\mathbb{R}.$$
\end{definition}
\end{frame}

\begin{frame}{Propiedades de la densidad}
\begin{theorem}[Propiedades de $f(x)$]
\begin{enumerate}
\item $f(x) \ge 0$ para todo $x$.
\item $\displaystyle\int_{-\infty}^{\infty} f(x)\,dx = 1$.
\item Donde $F$ es derivable, $F'(x) = f(x)$.
\item $\displaystyle P(a\le X\le b) = \int_a^b f(x)\,dx$ (área bajo la curva).
\end{enumerate}
\end{theorem}

\begin{alertblock}{Consecuencia importante}
Si $X$ es continua, $P(X=a) = \int_a^a f(x)\,dx = 0$ para todo $a$. Por eso
$$P(a\le X\le b) = P(a< X\le b) = P(a\le X< b) = P(a<X<b).$$
No importa incluir o no los extremos.
\end{alertblock}
\end{frame}

\begin{frame}{Interpretación de $f(x)$}
\begin{alertblock}{Cuidado: $f(x)$ NO es una probabilidad}
$f(x)$ puede ser mayor que 1. Lo que integra 1 es el \textbf{área total} bajo la curva, no los valores de $f$. $f(x)$ es una \textbf{densidad}: probabilidad "por unidad de longitud" cerca de $x$.
\end{alertblock}

\begin{block}{Esperanza y varianza (caso continuo)}
$$E(X) = \int_{-\infty}^\infty x\,f(x)\,dx, \qquad V(X) = \int_{-\infty}^\infty (x-E(X))^2 f(x)\,dx = E(X^2)-[E(X)]^2.$$
\end{block}

\begin{example}
Sea $f(x) = 2x$ para $0\le x\le 1$ (y 0 fuera). Verificamos: $\int_0^1 2x\,dx = [x^2]_0^1=1$. Entonces
$$P(0.3\le X\le 0.6) = \int_{0.3}^{0.6}2x\,dx = [x^2]_{0.3}^{0.6} = 0.36-0.09=0.27.$$
\end{example}
\end{frame}

\begin{frame}{Gráfico: densidad y área como probabilidad}
\begin{center}
\begin{tikzpicture}
\begin{axis}[
  xmin=0, xmax=1.1, ymin=0, ymax=2.2,
  xlabel={$x$}, ylabel={$f(x)$},
  width=10cm, height=5.5cm,
  axis lines=middle,
  domain=0:1, samples=60,
]
\addplot[thick,blue] {2*x};
\addplot[blue!20, domain=0.3:0.6] {2*x} \closedcycle;
\addplot[dashed] coordinates {(0.3,0)(0.3,0.6)};
\addplot[dashed] coordinates {(0.6,0)(0.6,1.2)};
\end{axis}
\end{tikzpicture}
\end{center}
El área sombreada entre $x=0.3$ y $x=0.6$ representa $P(0.3\le X\le 0.6)=0.27$.
\end{frame}

\section{Distribución Uniforme continua}

\begin{frame}{El modelo continuo más simple}
\begin{block}{Motivación}
Así como la Uniforme discreta modela ``todos los valores igualmente posibles'' en un conjunto finito, la \textbf{Uniforme continua} modela una v.a. que puede tomar \textbf{cualquier valor} en un intervalo $[a,b]$, sin que ningún subintervalo sea más probable que otro de igual longitud.
\end{block}

\begin{definition}[Distribución Uniforme continua]
$X\sim U(a,b)$ si su densidad es
$$f(x) = \begin{cases} \dfrac{1}{b-a}, & a\le x\le b, \\[4pt] 0, & \text{en el resto.}\end{cases}$$
\end{definition}
\end{frame}

\begin{frame}{Gráfico y función de distribución}
\begin{columns}
\begin{column}{0.48\textwidth}
\begin{center}
\begin{tikzpicture}
\begin{axis}[
  xmin=0, xmax=4.5, ymin=0, ymax=0.9,
  xlabel={$x$}, ylabel={$f(x)$},
  width=6cm, height=5cm,
  axis lines=left,
  xtick={1,3.5}, xticklabels={$a$,$b$}, ytick=\empty,
]
\addplot[thick,blue] coordinates {(0,0)(1,0)};
\addplot[thick,blue] coordinates {(1,0.6)(1,0)};
\addplot[thick,blue] coordinates {(1,0.6)(3.5,0.6)};
\addplot[thick,blue] coordinates {(3.5,0.6)(3.5,0)};
\addplot[thick,blue] coordinates {(3.5,0)(4.5,0)};
\end{axis}
\end{tikzpicture}
\end{center}
\end{column}
\begin{column}{0.5\textwidth}
\begin{block}{Función de distribución}
$$F(x) = \begin{cases} 0, & x<a, \\[2pt] \dfrac{x-a}{b-a}, & a\le x<b, \\[4pt] 1, & x\ge b.\end{cases}$$
\end{block}
\end{column}
\end{columns}
La densidad es constante en $[a,b]$ (rectángulo de altura $\frac{1}{b-a}$ y base $b-a$: área $1$); $F$ crece linealmente en ese tramo.
\end{frame}

\begin{frame}{Esperanza, varianza y función generadora de momentos}
\begin{theorem}
Si $X\sim U(a,b)$:
$$E(X) = \frac{a+b}{2}, \qquad V(X) = \frac{(b-a)^2}{12}, \qquad M_X(t) = \frac{e^{tb}-e^{ta}}{t(b-a)} \ (t\ne 0), \quad M_X(0)=1.$$
\end{theorem}

\begin{proof}[Cálculo de $E(X)$]
$$E(X) = \int_a^b x\,\frac{1}{b-a}\,dx = \frac{1}{b-a}\left[\frac{x^2}{2}\right]_a^b = \frac{b^2-a^2}{2(b-a)} = \frac{a+b}{2}. \qedhere$$
\end{proof}

\begin{alertblock}{Interpretación}
La media es el punto medio del intervalo (por simetría); la varianza depende solo de la longitud $b-a$: a mayor intervalo, mayor dispersión.
\end{alertblock}
\end{frame}

\begin{frame}{Ejemplo resuelto: Uniforme}
\begin{example}
Un colectivo pasa por una parada cada 15 minutos, y un pasajero llega en un instante aleatorio. Sea $X$ el tiempo de espera, $X\sim U(0,15)$. (a) Hallar $E(X)$ y $V(X)$. (b) ¿Cuál es la probabilidad de esperar más de 10 minutos?
\end{example}

\begin{block}{Resolución}
\textbf{(a)} $E(X) = \dfrac{0+15}{2} = 7.5$ min, \quad $V(X) = \dfrac{(15-0)^2}{12} = 18.75$ min$^2$.

\textbf{(b)} $P(X>10) = 1-F(10) = 1-\dfrac{10-0}{15-0} = \dfrac{5}{15} = 0.\overline{3}$.
\end{block}
\end{frame}

\section{La distribución Normal}

\begin{frame}{La distribución más importante de la estadística}
\begin{block}{Por qué es tan importante}
\begin{itemize}
\item Describe con buena aproximación numerosos fenómenos naturales y sociales (estaturas, errores de medición, pesos, calificaciones).
\item Es el límite de sumas de muchas variables aleatorias independientes (\textbf{Teorema Central del Límite}, Unidad 3).
\item Es la base de gran parte de la inferencia estadística (intervalos de confianza, tests de hipótesis).
\end{itemize}
\end{block}

{\footnotesize Fue descubierta y publicada por primera vez en 1733 por De Moivre; a ella llegaron también, de forma independiente, Laplace (1812) y Gauss (1809), en relación con la teoría de errores de observación astronómica y física. Por eso también se la llama distribución de Laplace-Gauss.}
\end{frame}

\begin{frame}{La distribución Normal: definición}
\begin{definition}[Distribución Normal]
$X\sim N(\mu,\sigma^2)$ si su densidad es
$$f(x) = \frac{1}{\sigma\sqrt{2\pi}}\exp\left(-\frac{(x-\mu)^2}{2\sigma^2}\right), \qquad x\in\mathbb{R},$$
con parámetros $\mu\in\mathbb{R}$ (media) y $\sigma>0$ (desvío estándar).
\end{definition}
\end{frame}

\begin{frame}{Propiedades de la campana de Gauss}
\begin{theorem}
Si $X\sim N(\mu,\sigma^2)$, entonces $E(X)=\mu$ y $V(X)=\sigma^2$.
\end{theorem}

\begin{block}{Propiedades geométricas de $f(x)$}
\begin{itemize}
\item Simétrica respecto de $x=\mu$ (forma de "campana").
\item Máximo en $x=\mu$, con $f(\mu) = \dfrac{1}{\sigma\sqrt{2\pi}}$.
\item Puntos de inflexión en $x=\mu\pm\sigma$.
\item $\mu$ controla el centro; $\sigma$ controla la dispersión (curva más "achatada" si $\sigma$ grande).
\end{itemize}
\end{block}

\begin{alertblock}{Regla empírica}
$P(\mu-\sigma<X<\mu+\sigma)\approx 0.68$; \quad $P(\mu-2\sigma<X<\mu+2\sigma)\approx 0.95$; \quad $P(\mu-3\sigma<X<\mu+3\sigma)\approx 0.997$.
\end{alertblock}
\end{frame}

\begin{frame}{Gráfico: efecto de $\mu$ y $\sigma$}
\begin{center}
\begin{tikzpicture}
\begin{axis}[
  xmin=-6, xmax=8, ymin=0, ymax=0.85,
  xlabel={$x$}, ylabel={$f(x)$},
  width=10.5cm, height=5.5cm,
  domain=-6:8, samples=100,
  legend pos=north east, legend style={font=\tiny},
]
\addplot[thick,blue] {1/(1*sqrt(2*pi))*exp(-(x-0)^2/(2*1^2))};
\addplot[thick,red] {1/(1*sqrt(2*pi))*exp(-(x-3)^2/(2*1^2))};
\addplot[thick,green!60!black] {1/(2*sqrt(2*pi))*exp(-(x-0)^2/(2*2^2))};
\legend{{$N(0,1)$}, {$N(3,1)$}, {$N(0,4)$}}
\end{axis}
\end{tikzpicture}
\end{center}
Desplazar $\mu$ traslada la curva; aumentar $\sigma$ la "achata" y ensancha (manteniendo área 1).
\end{frame}

\begin{frame}{Función generadora de momentos de la Normal}
\begin{theorem}
Si $X\sim N(\mu,\sigma^2)$:
$$M_X(t) = \exp\left(\mu t + \frac{\sigma^2 t^2}{2}\right), \qquad t\in\mathbb{R}.$$
\end{theorem}

\begin{block}{Idea de la deducción (caso $Z\sim N(0,1)$)}
$$M_Z(t) = \int_{-\infty}^\infty e^{tz}\frac{1}{\sqrt{2\pi}}e^{-z^2/2}\,dz = \int_{-\infty}^\infty \frac{1}{\sqrt{2\pi}}e^{-\frac{(z-t)^2}{2}+\frac{t^2}{2}}\,dz = e^{t^2/2},$$
completando cuadrados y usando que la última integral es el área total de una densidad $N(t,1)$, que vale $1$.
\end{block}
\end{frame}

\begin{frame}{Función generadora de momentos de la Normal (cont.)}
\begin{block}{De $Z$ a $X$}
Para $X=\mu+\sigma Z$ (ver estandarización, próxima sección):
$$M_X(t)=E(e^{t(\mu+\sigma Z)}) = e^{\mu t}M_Z(\sigma t) = e^{\mu t + \sigma^2t^2/2}.$$
\end{block}

\begin{alertblock}{Verificación rápida}
$M_X'(0)=\mu=E(X)$; puede verificarse que $M_X''(0)-[M_X'(0)]^2 = \sigma^2 = V(X)$.
\end{alertblock}
\end{frame}

\section{Estandarización}

\begin{frame}{La distribución Normal estándar}
\begin{block}{Problema}
$F(x)=\int_{-\infty}^x f(t)\,dt$ para la Normal \textbf{no tiene primitiva elemental}. Se resuelve numéricamente y se tabula para un único caso: la \textbf{Normal estándar}.
\end{block}

\begin{definition}[Normal estándar]
$Z\sim N(0,1)$ tiene densidad
$$\varphi(z) = \frac{1}{\sqrt{2\pi}}e^{-z^2/2},$$
y función de distribución $\Phi(z) = P(Z\le z)$, tabulada.
\end{definition}
\end{frame}

\begin{frame}{La distribución Normal estándar: estandarización}
\begin{theorem}[Estandarización]
Si $X\sim N(\mu,\sigma^2)$, entonces
$$Z = \frac{X-\mu}{\sigma} \sim N(0,1).$$
\end{theorem}
Esto permite calcular probabilidades de \textbf{cualquier} Normal usando la tabla de $\Phi$.
\end{frame}

\begin{frame}{Uso de la tabla y simetría de $\Phi$}
\begin{block}{Propiedades útiles de $\Phi$}
\begin{itemize}
\item $\Phi(0) = 0.5$.
\item $\Phi(-z) = 1-\Phi(z)$ (simetría de la densidad respecto de 0).
\item $P(Z>z) = 1-\Phi(z)$.
\item $P(a<Z<b) = \Phi(b)-\Phi(a)$.
\end{itemize}
\end{block}

\begin{example}
Con tabla estándar: $\Phi(1.96)\approx 0.975$, entonces $P(Z<1.96)=0.975$ y $P(Z>1.96)=0.025$.

Por simetría: $\Phi(-1.96) = 1-0.975 = 0.025$.
\end{example}
\end{frame}

\begin{frame}{Ejemplo resuelto: Normal}
\begin{example}
El peso de los paquetes producidos por una máquina se distribuye $N(\mu=500\text{ g}, \sigma=15\text{ g})$. (a) ¿Qué proporción de paquetes pesa menos de 480 g? (b) ¿Entre 490 y 520 g?
\end{example}

\begin{block}{Resolución}
\textbf{(a)} $Z = \dfrac{480-500}{15} = -1.33$.
$$P(X<480) = P(Z<-1.33) = \Phi(-1.33) = 1-\Phi(1.33) \approx 1-0.9082 = 0.0918.$$

\textbf{(b)} $Z_1 = \dfrac{490-500}{15}=-0.67$, \quad $Z_2 = \dfrac{520-500}{15}=1.33$.
$$P(490<X<520) = \Phi(1.33)-\Phi(-0.67) = 0.9082 - (1-0.7486) = 0.9082-0.2514 = 0.6568.$$
\end{block}
\end{frame}

\begin{frame}{Ejemplo resuelto: percentiles (problema inverso)}
\begin{example}
Con los mismos datos ($\mu=500$, $\sigma=15$), ¿cuál es el peso $x_0$ tal que solo el 10\% de los paquetes pesa más que $x_0$?
\end{example}

\begin{block}{Resolución}
Buscamos $x_0$ tal que $P(X>x_0)=0.10$, es decir $\Phi(z_0)=0.90$. De la tabla, $z_0\approx 1.28$.

Deshaciendo la estandarización:
$$z_0 = \frac{x_0-\mu}{\sigma} \;\Rightarrow\; x_0 = \mu + z_0\sigma = 500 + 1.28\times 15 = 519.2 \text{ g}.$$
\end{block}
\end{frame}

\section{Resumen}

\begin{frame}{Resumen de la clase}
\begin{block}{Ideas clave}
\begin{itemize}
\item $X$ continua $\Leftrightarrow$ existe densidad $f\ge 0$ con $F(x)=\int_{-\infty}^x f(t)dt$ y $\int_{-\infty}^\infty f=1$.
\item $P(X=a)=0$ para toda v.a. continua; las probabilidades son áreas bajo $f$.
\item $X\sim U(a,b)$: densidad constante $1/(b-a)$; $E(X)=(a+b)/2$, $V(X)=(b-a)^2/12$.
\item $X\sim N(\mu,\sigma^2)$: simétrica, campana centrada en $\mu$, dispersión $\sigma$; $E(X)=\mu$, $V(X)=\sigma^2$, $M_X(t)=e^{\mu t+\sigma^2t^2/2}$.
\item Estandarización $Z=(X-\mu)/\sigma \sim N(0,1)$ permite usar la tabla de $\Phi$ para cualquier Normal.
\end{itemize}
\end{block}

\begin{alertblock}{Próxima clase}
Estudiaremos las distribuciones Exponencial y Gamma, y cómo hallar la distribución de una función $Y=g(X)$ de una variable aleatoria.
\end{alertblock}
\end{frame}

\end{document}
